\section{El progreso de la generación de video no es solo «una imagen más nítida»}

Esta clase comienza con varios videos generados, pero la pregunta de fondo no es elegir cuál fotograma es el más bello, sino entender por qué entre 2022 y 2024 progresaron al mismo tiempo el movimiento, los reflejos, el lenguaje de cámara, el apego al texto y la capacidad de edición. El hilo conductor que propone el ponente es sobrio: comprimir el medio en una secuencia de tokens manejable, entrenar con Flow Matching un Transformer grande y unificado, y hacer crecer en conjunto los datos, el cómputo y el tamaño del modelo. Más adelante se desglosa esta frase en representación, objetivo de aprendizaje, arquitectura, datos y sistema de evaluación, cada uno cuantificable.

\subsection{Qué demuestra cada uno de los tres ejemplos iniciales}

Los ejemplos iniciales cumplen tres papeles distintos como evidencia: el koala que surfea muestra un sujeto complejo en movimiento, el fantasma en el espejo muestra coherencia geométrica y de reflejos, y la edición del corredor muestra la conservación del contenido del video de entrada y el control mediante instrucciones. Esta sección establece primero un marco de lectura de figuras: primero preguntar «cuál es la tarea», luego «qué propiedades se conservaron» y por último «si se trata de una capacidad estable»; de lo contrario, unos cuantos fragmentos exitosos se presentan fácilmente, por error, como confiabilidad completa, y tampoco aportan evidencia para los juicios sobre la arquitectura que vienen después.

\lecturefigure{slide-064-moviegen-demo-koala.jpg}{Ejemplo de texto a video: el sujeto, la cámara, las olas y el movimiento de surf deben mantenerse en conjunto}{00:01:50}

La dificultad del fragmento de surf no es generar un koala, sino lograr que la apariencia del sujeto, la tabla, las olas, la distancia de la cámara y el movimiento continuo sean coherentes a lo largo de varios fotogramas. Un modelo de imagen de un solo fotograma solo necesita optimizar la plausibility espacial; un modelo de video, además, debe controlar en el eje temporal la estabilidad de la identidad, la continuidad del movimiento y la recuperación tras oclusiones.

\lecturefigure{slide-091-moviegen-demo-ghost.jpg}{Ejemplo del espejo: el reflejo, el punto de vista y la identidad del sujeto forman restricciones entre regiones}{00:02:22}

La escena del espejo exige que el modelo vincule el «sujeto real» y el «sujeto en el espejo» como un mismo objeto, respetando a la vez los cambios de punto de vista. Esto respalda que el modelo ya aprendió ciertas regularidades visuales, pero no demuestra que tenga un motor óptico explícito; una prueba más confiable requiere cambiar el ángulo del espejo, las oclusiones y la posición del sujeto, y observar si la regularidad se generaliza de manera sistemática.

\lecturefigure{slide-112-moviegen-editing-intro.jpg}{Edición de video por instrucciones: sobre una misma entrada se agrega un objeto, se cambia la ropa y se modifica el entorno}{00:02:48}

Las dimensiones de evaluación de la tarea de edición difieren de las de la generación pura. Además de si la salida se ve bien, hay que verificar si las regiones no editadas se conservan, si la región objetivo cumple estrictamente la instrucción, si el ritmo del movimiento continúa y si la modificación persiste a lo largo de todo el video. Más adelante se verá que este tipo de capacidad no surge automáticamente del modelo T2V base, sino que requiere una ruta de entrenamiento adicional.

\begin{warningbox}{Un ejemplo exitoso no es confiabilidad a nivel de distribución}
Una demostración en video puede probar que «existe al menos una salida exitosa», pero no que el modelo sea estable frente a la distribución de prompts, las semillas aleatorias, los objetos de cola larga y las interacciones complejas. Una evaluación confiable requiere un conjunto fijo de prompts, muestras de fallos públicas, varias semillas aleatorias y comparaciones humanas por pares; si solo se muestran los mejores ejemplos, no es posible estimar la tasa de éxito.
\end{warningbox}

\teachervoice{00:01:13--00:03:42, el docente usa tres demos iniciales para distinguir calidad de generación, reflejos complejos y edición por instrucciones. Enfatiza que la mejora de los últimos dos años no es la estética de un solo fotograma, sino que los modelos empezaron a manejar movimiento, relaciones y control.}

\subsection{De Make-A-Video a Movie Gen: qué cambió realmente}

Las dos figuras siguientes reúnen la línea de tiempo y las conclusiones de la clase. La primera compara las diferencias visibles entre 2022 y 2024; la segunda atribuye la mejora a una arquitectura unificada y al escalamiento, no a algún trick aislado. Para entender esta historia, conviene considerar los algoritmos, los datos, el hardware, la infraestructura de entrenamiento y los métodos de evaluación como variables conjuntas.

\lecturefigure{slide-150-video-generation-progress.jpg}{Comparación directa del progreso en dos años: Make-A-Video y Movie Gen}{00:03:35}

El modelo temprano de la izquierda ya podía generar contenido dinámico, pero los bordes del sujeto, la estabilidad del movimiento y la continuidad de los detalles son claramente más débiles; Movie Gen, a la derecha, es más sólido en las olas complejas, el seguimiento de cámara y la coherencia del sujeto. La figura respalda que «la frontera de capacidades se desplazó de forma notable», pero a partir de la sola comparación visual no es posible separar qué parte de la ganancia proviene de la arquitectura y qué parte de la escala de entrenamiento.

\lecturefigure{slide-155-moviegen-paper-thesis.jpg}{Tesis central de la clase: un Transformer sencillo con datos, cómputo y parámetros escalados también funciona para video}{00:05:02}

Aquí «sencillo» es relativo a las cascadas de U-Net y a los módulos de video altamente especializados, y no significa que el sistema sea barato. Movie Gen sigue necesitando TAE, varios codificadores de texto, Flow Matching, paralelismo distribuido, limpieza de datos, entrenamiento progresivo y postentrenamiento. La simplificación central ocurre en el backbone: reutilizar en lo posible un Transformer cuyo escalamiento es predecible, en lugar de reinventar la red para cada resolución y modalidad.

\begin{importantbox}{Implicaciones de ingeniería de unificar la arquitectura}
Unificar la arquitectura no significa que «todas las modalidades sean idénticas», sino concentrar la complejidad propia de cada modalidad en el encoder, el decoder y la conditioning interface, de modo que el tronco siga usando Transformer blocks escalables. Así se pueden reutilizar el modelo, los kernels, las estrategias de paralelismo y la infraestructura de entrenamiento, y la dirección de escalamiento también queda más clara.
\end{importantbox}

\subsection{Trayectoria histórica: de modelos pequeños especializados a modelos grandes unificados}

Los ejemplos anteriores mostraban resultados; esta sección pasa a la historia de los mecanismos. La línea de tiempo debe responder dos preguntas: en qué estructuras especializadas se apoyaban los métodos tempranos y por qué los modelos grandes de 2024 prefieren un Transformer unificado. Al leer la figura, no hay que equiparar automáticamente «nuevo» con «mejor»; lo decisivo es qué sesgo se beneficia de manera más sostenida una vez que cambian la escala de entrenamiento y la infraestructura.

\lecturefigure{slide-174-video-generation-history.jpg}{Trayectoria de la generación de video: difusión, cascadas de U-Net y unificación con Transformer/Flow Matching}{00:09:39}

Los primeros Video Diffusion, Make-A-Video e Imagen Video solían combinar U-Net, difusión, interpolación de fotogramas y superresolución en varias etapas dentro de un pipeline complejo. Con líneas como Sora y Movie Gen, el foco de la investigación se desplazó hacia un backbone unificado de tokens, Flow Matching y entrenamiento a mayor escala. El cambio no es que todos los módulos antiguos desaparezcan, sino que el sistema empieza a dar prioridad a reducir las bifurcaciones de la arquitectura.

\lecturefigure{slide-175-moviegen-overview.jpg}{Panorama de Movie Gen Video: escala, tareas y datos de entrenamiento}{00:10:37}

Movie Gen Video es un joint text-to-image/text-to-video foundation model de unos 30B parámetros, con un contexto máximo de unos 73K video tokens, que corresponde a una configuración de entrenamiento con videos de hasta 16 segundos a 16 fps. El ponente describe el volumen de datos con órdenes de magnitud de $O(100\mathrm{M})$ videos y $O(1\mathrm{B})$ imágenes; son cifras de escala y no deben interpretarse como conteos exactos publicados.

\begin{knowledgebox}{Por qué entrenar imágenes y video de forma conjunta}
Los datos de imagen son más abundantes y tienen una cobertura semántica más amplia, lo que ayuda al modelo a aprender objetos, estilos y composición; los datos de video aportan movimiento y relaciones temporales. El entrenamiento conjunto permite que un mismo backbone absorba a la vez regularidades espaciales y espaciotemporales, y también que una imagen entre en la representación unificada como un video de un solo fotograma.
\end{knowledgebox}

\teachervoice{00:07:57--00:11:19, el docente resume la trayectoria como el paso de «arquitecturas especializadas y de pequeña escala» a «arquitectura unificada, Transformer sencillo y escalamiento». Al mismo tiempo da la workload real de 30B, 73K tokens, 16 segundos y 16 fps, y advierte que unificar no equivale a abaratar.}

\subsection{Resumen de la sección}

El progreso de la generación de video abarca calidad espacial, coherencia temporal, relaciones complejas y edición controlable. La tesis central de Movie Gen es unificar el backbone con un Transformer escalable, pero el sistema completo sigue dependiendo de la compresión de la representación, el objetivo de aprendizaje, los datos, el paralelismo y la evaluación. Los ejemplos son evidencia de capacidad, no prueba suficiente de confiabilidad ni de un modelo causal del mundo.
